\section{Existence, least periods and oriented representatives}
\label{sec:least}

For each row of Table~\ref{tab:classification}, substitute its
displayed entries and parameter into
$x_{i+3}+x_i-x_{i+1}x_{i+2}-a$. The result is identically
zero at every wraparound phase. This is a polynomial identity
in the displayed free integer variables, so it establishes
existence on the full parameter range, not just on the finite
core. Substitution of the initial triple in $K_a$ gives exactly
the displayed level. Invariance then gives that same value at
all other phases. The preserved separate symbolic check verifies
these recurrence and level identities at all $55$ phases across
the ten rows, in addition to the full invariant identity;
the identities themselves are immediate from the printed words.

We give all period and identification arguments, since a periodic
word alone is not a disjoint list of native cycles.
For $F_1$ and $F_2$, Lemma~\ref{lem:zero} already gives
the least period, and $u<v$ chooses the unique ordering of an
alternating pair. The length-three word $(-2,-2,t)$ is constant
only at $t=-2$. Thus $t\ne-2$ gives exactly period three;
its parameter is uniquely recovered from $a=2t-4$.

For $F_4$, shifting by two positions exchanges $t$ and
$a+1-t$. A shorter period is possible exactly when
$2t=a+1$, giving an alternating word. At fixed $a$, the
strict inequality $2t<a+1$ removes that case and chooses
one representative of each exchange. Any further cyclic
identification must preserve the alternating positions carrying
$-1$; if one of the free entries also equals $-1$, checking
the four positions gives the same exchange and no extra
identification.

The $F_5$ word has prime length five and is never constant,
as it contains both $0$ and $-1$. For
$t\notin\{0,-1\}$ its displayed occurrence of $-1$ is
unique. A cyclic identification of two such words must align
that occurrence, forcing equal parameters. The words at $t=0$
and $t=-1$ are rotations of one another and are the sole
duplication. Omitting $t=-1$ therefore gives the stated range.
In particular, parameters $t$ and $-1-t$ generally give
different oriented cycles, exchanged by time reversal. The
exceptional five-cycle $E_5$ is nonconstant and lies at
the different parameter $a=-13$.

For $F_6$, the condition $t\ne0$ excludes every proper
divisor $1,2,3$ of six. A shift by three positions exchanges
$t$ with $-t$, and the two nonzero positions show that these
are the only possible parameter identifications. The choice
$t\ge1$ is consequently exact.

For $F_8$, a shift by four positions exchanges $t$ with
$-t-2$. Period four would force $t=-1$, and periods one
or two would imply period four. The excluded value $t=-1$
is already $F_4$ at level $1$. For $t\ge0$, the positions
carrying $-1$ occur precisely four apart, so a cyclic
identification can only give the stated exchange. Exactly one
member of each nondegenerate pair has $t\ge0$.

For $F_{12}$ with $m\ge2$, compare entries with zero-based
indices. Periods one or two are impossible because the entries
at positions $0$ and $4$ are $1$ and $-1$.
Period three is excluded by positions $1$ and $4$; period
four by positions $0$ and $4$; period six by positions $1$
and $7$, which are $m$ and $1-m$.
For $m=1$ the word is
\[
 (1,1,1,0,-1,-1,1,0,1,-1,-1,0).
\]
Positions $0,3$ exclude period three, positions $0,4$
exclude period four, and positions $1,7$ exclude period six.
Periods one and two are again excluded by positions $0,4$.
These are every proper divisor of twelve. Since
$m^2-m+2$ is strictly increasing for integers $m\ge1$,
different standard parameters cannot describe the same cycle.
This is a direct least-period argument; no unprinted appeal
to the unforced predecessor is required.

The length-nine exceptional word has different entries in
positions $0$ and $3$, so it has neither period three nor
period one. It therefore has least period nine. The reversed
$E_5$ word is its rotation by one position, and the reversed
$E_9$ word is its rotation by six. Thus no additional
reverse-oriented exceptions are hidden by the list.

Finally, rows of different least periods cannot overlap, and
the two rows of length five have different forcing parameters.
Together with Section~\ref{sec:finite}, these observations
complete every assertion of Theorem~\ref{thm:classification}.
